\scanpage{53}
\begin{flushright}
at\illegible{1} \mspage{13}
\end{flushright}

\noindent exist for the smaller primes as well. As an example with $q=5$ we find:

\uline{Let}*$t(n)=c(5^{\alpha}n)/c(5^{\alpha})$. \uline{Then}
\par\noindent\rule{\linewidth}{0.3pt}\par
\noindent * Since\reversemarginpar\margin{\handwritten{on line}} $\handwritten{\mathord{\tikz[x=1ex,y=1ex,baseline=0pt]{\draw (.82,1.85) .. controls (.42,1.55) and (.34,-.54) .. (0,-.35);}}}\raisebox{.37ex}{$5^{\alpha+1}$}\handwritten{\mathord{\tikz[x=1ex,y=1ex,baseline=0pt]{\draw (0,1.62) .. controls (.04,1.02) and (-.13,.29) .. (.35,.1);}}}\handwritten{\mid}c(5^{\alpha}n)$ and $5^{\alpha+1}\handwritten{\parallel}c(5^{\alpha})$ this definition is legitimate in any residue class ring modulo $5^{\beta}$.
\par\noindent\rule{\linewidth}{0.3pt}\par
\[
\begin{aligned}[b]
&t(5n)\handwritten{\equiv}t(5)t(n)\qquad(\mathrm{mod}\;5^{3\alpha}),\\
\text{\uline{and} }\;&t(np)-t(n)t(p)+p^{-1}t(n/p)\handwritten{\equiv}0\quad(\mathrm{mod}\;5^{3\alpha})\quad\text{\uline{if} }p\mathrel{\rlap{$\handwritten{/}$}=}5\ \text{\uline{is} \uline{prime}.}
\end{aligned}
\tag*{(7.1)}
\]

For $q=2,3,5,\text{and }7$ these results can be fully proved as for $q=13$; for $q=11$ we have again the same reservations as for $q=19$. In a s\makebox[0pt][l]{\handwritten{\tikz[x=1em,y=1ex,baseline=0pt]{\draw (.02,-.28)--(.54,-.28)--(.54,1.54)--(1.01,1.6);}}}nese we may feel that the Ramanujan-type properties (2.5) which do not generalise to higher primes have previously concealed the existence of the multiplicative properties. Thus(to suppose what is in fact false) had $c(2^{\alpha}\inserted{\handwritten{\text{\uline{$n$}}}})$ been prime to $2$,its multiplicative properties might have been sought and found by hand; as it is\handwritten{,}the first non-Ramanujan prime is $13$ which is beyond the scope of hand computation. On the other hand had $p(5n+4)$ been prime to $5$ the whole possibility of considering congruences for modular forms m\inserted{\handwritten{i}}ght never have been noticed.

We may finally try to summarise the distinction between Ramanujan-type and multiplicative properties in this way. The Ramanujan-type property asserts that
\[
\handwritten{\sum_{n=1}^{\infty}}c(5^{\alpha}n)\inserted{\handwritten{x^{n}}}=5^{\alpha+1}P(x),
\]
where \uline{$P(x$}$)$ is a power series with integral coefficients. The multiplicative property now asserts that \uline{$P(x$}$)$ itself is an interesting series.

8. \uline{Notes} \uline{on} \uline{computation}. The main requirement for most of this work was the ability to compute the Fo\makebox[0pt][l]{\handwritten{\tikz[x=1em,y=1ex,baseline=0pt]{\draw (.02,-.28)--(.54,-.28)--(.54,1.54)--(1.01,1.6);}}}ruier expansions of various modular forms and functions modulo \uline{$m$}. If we have available in store various power series $g_{i}(x)$ up to the term in $x^{N}$,then it is clear that addition and subtraction of constant multiples of these series,testing for ``accidents'',etc.,takes a time approximately proportional to $N$; we shall abbreviate this to ``takes $O(N)$''. The test for M-results takes
\normalmarginpar
% archive sheet number: 52
